\documentclass[12pt]{article}
\usepackage[utf8]{inputenc}
\usepackage[T1]{fontenc}
\usepackage{hyperref}
\title{Serili Language Resources for K-12 Curriculum Development}
\author{Compiled by Claude Sonnet 5.5 (claude-sonnet-5-5) for Andr\'as Kornai}
\date{}
\begin{document}
\maketitle

\section{Introduction}
Serili (ISO 639-3: sve; Glottocode seri1255) is an Austronesian (Malayo-Polynesian, Babar group) language spoken on Masela Island in South Maluku, Indonesia. The figure of about 330 speakers dates from 1980. The language has no established written tradition. The search found only catalogue and database entries.

\section{Primary Educational Resources}
No primers, textbooks or school programmes for Serili were found.

\section{Linguistic and Reference Materials}
\begin{itemize}
\item Wikipedia, Serili language: \url{https://en.wikipedia.org/wiki/Serili_language}
\item Verbix endangered-language page for Serili: \url{https://docs.verbix.com/EndangeredLanguages/Serili}
\item ASJP wordlist for Serili: \url{https://asjp.clld.org/languages/SERILI}
\item Endangered Languages Project entry: \url{https://endangeredlanguages.com/lang/2781}
\end{itemize}

\section{Religious and Community Resources}
\begin{itemize}
\item Joshua Project, Serili people profile (about 600 people, mostly Christian): \url{https://joshuaproject.net/people_groups/14870}
\item Joshua Project language page for sve: \url{https://joshuaproject.net/languages/sve}
\end{itemize}

\section{Digital and Audio Resources}
No audio, video or corpus resources beyond the ASJP wordlist were found.

\section{Curriculum Development Context}
Almost nothing is available for Serili. The language is unwritten, the speaker data is old, and no school use is known. Any curriculum would require new documentation work with the community, and Indonesian would remain the language of instruction in the meantime.

\end{document}